% Official source: https://www.imo-official.org/assets/documents/problems/2022/2022_eng.pdf
\competitionyear{2022}

\begin{problem}{2022}{1}
The Bank of Oslo issues two types of coin: aluminium (denoted $A$) and bronze (denoted $B$). Marianne has $n$ aluminium coins and $n$ bronze coins, arranged in a row in some arbitrary initial order. A \emph{chain} is any subsequence of consecutive coins of the same type. Given a fixed positive integer $k\le2n$, Marianne repeatedly performs the following operation: she identifies the longest chain containing the $k$-th coin from the left, and moves all coins in that chain to the left end of the row. For example, if $n=4$ and $k=4$, the process starting from the ordering $AABBBABA$ would be
\[
\begin{gathered}
    AAB\underline{B}BABA\to BBB\underline{A}AABA\to AAA\underline{B}BBBA\\
    \to BBB\underline{B}AAAA\to BBB\underline{B}AAAA\to\cdots.
\end{gathered}
\]
Find all pairs $(n,k)$ with $1\le k\le2n$ such that for every initial ordering, at some moment during the process, the leftmost $n$ coins will all be of the same type.
\end{problem}

\begin{problem}{2022}{2}
Let $\R_{>0}$ denote the set of positive real numbers. Find all functions $f\colon\R_{>0}\to\R_{>0}$ such that for each $x\in\R_{>0}$, there is exactly one $y\in\R_{>0}$ satisfying
\[
    xf(y)+yf(x)\le2.
\]
\end{problem}

\begin{problem}{2022}{3}
Let $k$ be a positive integer and let $S$ be a finite set of odd prime numbers. Prove that there is at most one way (up to rotation and reflection) to place the elements of $S$ around a circle such that the product of any two neighbours is of the form $x^2+x+k$ for some positive integer $x$.
\end{problem}

\begin{problem}{2022}{4}
Let $ABCDE$ be a convex pentagon such that $BC=DE$. Assume that there is a point $T$ inside $ABCDE$ with $TB=TD$, $TC=TE$ and $\angle ABT=\angle TEA$. Let line $AB$ intersect lines $CD$ and $CT$ at points $P$ and $Q$, respectively. Assume that the points $P,B,A,Q$ occur on their line in that order. Let line $AE$ intersect lines $CD$ and $DT$ at points $R$ and $S$, respectively. Assume that the points $R,E,A,S$ occur on their line in that order. Prove that the points $P,S,Q,R$ lie on a circle.
\end{problem}

\begin{problem}{2022}{5}
Find all triples $(a,b,p)$ of positive integers with $p$ prime and
\[
    a^p=b!+p.
\]
\end{problem}

\begin{problem}{2022}{6}
Let $n$ be a positive integer. A \emph{Nordic square} is an $n\times n$ board containing all the integers from $1$ to $n^2$ so that each cell contains exactly one number. Two different cells are considered adjacent if they share a common side. Every cell that is adjacent only to cells containing larger numbers is called a \emph{valley}. An \emph{uphill path} is a sequence of one or more cells such that:
\begin{problemclauses}
    \item the first cell in the sequence is a valley,
    \item each subsequent cell in the sequence is adjacent to the previous cell, and
    \item the numbers written in the cells in the sequence are in increasing order.
\end{problemclauses}
Find, as a function of $n$, the smallest possible total number of uphill paths in a Nordic square.
\end{problem}
